% Lec 04：ZFC、算术化与 Gödel 不完备定理
\documentclass[UTF8,a4paper,11pt]{ctexart}
\usepackage[a4paper,margin=2.3cm]{geometry}
\usepackage{amsmath,amssymb,mathtools}
\usepackage{enumitem}
\usepackage[hidelinks]{hyperref}
\usepackage{fancyhdr}
\setlist[itemize]{leftmargin=2.2em,itemsep=0.25em,topsep=0.35em}
\setlength{\parindent}{2em}
\setlength{\parskip}{0.35em}
\pagestyle{fancy}\fancyhf{}
\fancyhead[L]{离散数学与结构（I）}\fancyhead[R]{Lec 04}
\fancyfoot[C]{\thepage}\setlength{\headheight}{14pt}
\title{ZFC、算术化与 Gödel 不完备定理}
\author{zhangjunshuo}\date{2026 年 9 月}
\begin{document}
\maketitle
\begin{quote}
本讲从 Russell 悖论说明公理化集合论的必要性，再把公式和证明编码为自然数，最后由不动点引理连接停机问题与 Gödel 不完备定理。
\end{quote}
\tableofcontents\clearpage
\section{Russell 悖论与 ZFC}


\subsection{从朴素概括到 Russell 悖论}


若允许对任意性质直接造集合，就可以写

\[
R=\{x:x\notin x\}.
\]

于是 $R\in R$ 当且仅当 $R\notin R$，矛盾。ZFC 的出发点不是把“所有性质的外延”都当作集合，而是规定哪些集合构造由公理保证。

\subsection{ZFC 的九类公理}


1. \textbf{外延}：具有相同元素的集合相等；
2. \textbf{分离模式}：从已有集合 $a$ 中按公式筛选子集；
3. \textbf{配对}：由 $x,y$ 构造 $\{x,y\}$；
4. \textbf{并集}：由 $a$ 构造 $\bigcup a$；
5. \textbf{幂集}：由 $a$ 构造 $\mathcal P(a)$；
6. \textbf{无穷}：存在含 $\varnothing$ 且对后继封闭的集合；
7. \textbf{替换模式}：函数式关系在集合上的像仍是集合；
8. \textbf{正则}：每个非空集合有 $\in$-极小元；
9. \textbf{选择公理}：非空集合族有选择函数。

分离和替换是模式，公式中使用的参数可以自由出现，但新产生的集合变量不能在公式中自由出现。正则公理排除 $x\in x$、有限循环和无穷下降的 $\in$ 链。

\subsection{ZFC 的对象观}


自然数可由 $0=\varnothing$ 和后继 $n^+=n\cup\{n\}$ 构造；Kuratowski 有序对可写为

\[
(x,y)=\{\{x\},\{x,y\}\}.
\]

有序对、笛卡尔积、关系和函数都必须说明它们如何由已有集合经幂集、替换和分离得到。分离只筛选已有集合的子集，不能单独凭空造出“所有满足性质的对象”。
\clearpage
\section{语法与证明的算术化}


\subsection{编码方案只是示意}


给基本符号分配自然数，再用素数幂、进制或其他可计算编码把有限字符串编码成自然数。记公式 $\varphi$ 的编码为 $\ulcorner\varphi\urcorner$。具体常数和编码基数可以改变，不影响后续论证。

\subsection{最小闭包定义}


TERM、FORM 和 PROOF 都可以通过“最小闭包”定义：

\begin{itemize}
\item TERM 包含变量、常元，并对函数符号应用封闭；
\item FORM 包含原子式，并对联结词和量词应用封闭；
\item PROOF 包含公理行，并对每条推理规则应用封闭。
\end{itemize}


“最小”表示：它本身满足生成规则，且任何满足规则的集合都包含它。对自然数上的编码集合，这些定义给出可机械检查的句法。

\subsection{可证性谓词}


对有效公理化理论，存在原始递归关系

\[
\operatorname{Proof}(d,y)
\]

表示 $d$ 是以编码 $y$ 的句子为结论的有限合法证明。于是

\[
\operatorname{Prov}(y)\;\Longleftrightarrow\;\exists d\,\operatorname{Proof}(d,y).
\]

若公理集合只有可枚举而非可判定的给法，需要把公理出现的枚举阶段一并编码，才能保持证明检查可计算。不能只由“公式可以归纳定义”推出任意理论的定理集可枚举。

\subsection{代入函数}


设 $S(a,b)$ 把编码为 $a$ 的公式中指定自由变量替换为数词 $\overline b$，并返回新编码。$S$ 是后面不动点构造的自代入工具；它的正确性依赖变量捕获检查和有效编码。
\clearpage
\section{不动点引理与停机问题}


\subsection{不动点引理}


对只有一个自由变量的公式 $\varphi(x)$，存在句子 $\psi$ 使

\[
\mathrm{PA}\vdash\psi\leftrightarrow\varphi(\ulcorner\psi\urcorner).
\]

构造思路是令

\[
\theta(x):=\varphi(S(x,x)),
\]

取 $m=\ulcorner\theta\urcorner$，再令 $\psi=\theta(\overline m)$。此时 $S(m,m)$ 正是 $\psi$ 的编码，因而 $\psi$ 可以谈论自身的编码。

令 $\varphi(x)=\neg\operatorname{Prov}(x)$，得到 Gödel 句

\[
\psi\leftrightarrow\neg\operatorname{Prov}(\ulcorner\psi\urcorner).
\]

\subsection{停机问题}


不存在总程序 $H$，对每个程序编码 $\ulcorner M\urcorner$ 和输入 $x$ 判断 $M(x)$ 是否停机。若存在，构造程序 $D$：

1. 运行 $H(\ulcorner D\urcorner,\ulcorner D\urcorner)$；
2. 若 $H$ 说停机，则进入死循环；
3. 若 $H$ 说不停机，则立即停机。

运行 $D$ 自身会产生矛盾：$H$ 的预测永远与实际相反。

\subsection{与证明枚举的关系}


在有效公理化理论中，有限证明可以机械检查，所以定理集至少是递归可枚举的。若一个足够强的理论同时能判定每个句子或其否定，就可以并行枚举两类证明得到判定器；这会与相应的不可判定性结果冲突。这里的“证明可检查”和“理论足够强”都是必要前提。
\clearpage
\section{Gödel 不完备定理}


\subsection{第一不完备定理的形式}


设 $T$ 是：

\begin{itemize}
\item 一致的；
\item 有效公理化的；
\item 足够强，能够表达并验证基本算术的编码事实。
\end{itemize}


由不动点引理得到 Gödel 句 $G_T$，满足

\[
T\vdash G_T\leftrightarrow\neg\operatorname{Prov}_T(\ulcorner G_T\urcorner).
\]

在适当的可靠性条件下，$G_T$ 在标准自然数中为真但 $T$ 不能证明它。要排除 $\neg G_T$ 的可证性，可以假设 $\omega$-一致性，或改用 Rosser 句在较弱的一致性假设下得到独立性。

因此，结论不是“任何含算术的理论都不完备”，而是带有一致性、有效公理化和算术表达能力等条件的定理。

\subsection{第二不完备定理（概要）}


若 $T$ 还足够强，能在自身内部形式化“存在一个合法证明”的谓词和 Hilbert–Bernays 导出条件，则它不能在自身内部证明自己的形式一致性句子 $\mathrm{Con}(T)$，即

\[
T\nvdash\mathrm{Con}(T)
\]

（假设 $T$ 一致）。这里的 $\mathrm{Con}(T)$ 是依赖编码方案写出的算术句子，不是元语言中一句未经形式化的“我没有矛盾”。

\subsection{不要混淆三层陈述}


1. 元理论中说“某个有限字符串是合法证明”；
2. 算术理论中用 $\operatorname{Proof}$ 和 $\operatorname{Prov}$ 表示这件事；
3. 模型中讨论这些算术句子的真值。

Gödel 论证正是把三层通过编码连接起来；省略其中的有效性或一致性条件，就会把定理误说成无条件的句法悖论。

\end{document}
